% GENERATED FILE. Edit canonical lessons, then run npm run generate:books.
% canonical-source-hash: fnv1a64:e5063f3718d072f7
% canonical-lessons: TA-C234-ancalakam, TA-C234-nakaram, TA-C234-natu, TA-C234-talainakaram, TA-C234-narcanti

\chapter{Post office, City, Country, Capital city, Crossroads}
\label{ch:ta-post-office-city-country-capital-city-crossroads}
\hlchaptermodality{234}

\begin{chapteropening}
\textbf{By the end of this chapter:} \emph{I can say a post office, a city, a country, a capital city and a crossroads.}

Complete the last lesson of chapter 234: I can say a post office, a city, a country, a capital city and a crossroads.
\end{chapteropening}

\section[\texorpdfstring{\ta{அஞ்சலகம்} (añcalakam)}{அஞ்சலகம் (añcalakam)}]{\ta{அஞ்சலகம்} (añcalakam) — a post office}

\label{lesson:TA-C234-ancalakam}



\begin{quote}
Before the new one: say the Tamil for the day before yesterday, then the Tamil for a second (of time).
\end{quote}

\subsection*{You'll want to know: \ta{அஞ்சலகம்}}
\textbf{\ta{அஞ்சலகம்}} — \emph{añcalakam} — \textquotedblleft{}a post office\textquotedblright{}.

\begin{cousinweb}
\textbf{\ta{அஞ்சலகம்}} (\emph{añcalakam}) is native: \emph{añcal}, \textquotedblleft{}relay, letter post\textquotedblright{} (DEDR 54, which lists Malayalam añcal and Telugu ance), and \emph{akam}, \textquotedblleft{}house\textquotedblright{} (DEDR 7): \textquotedblleft{}the post house\textquotedblright{}.
\end{cousinweb}

\begin{grammarlens}[title={What you've built}]
One more word for this chapter.
\end{grammarlens}

\subsection*{Your turn}
\begin{itemize}
\raggedright
  \item \emph{Say it:} \emph{añcalakam}
  \item \emph{Say it:} \emph{añcalakam}, once more
  \item \emph{Say it:} say \emph{añcalakam}
  \item \emph{Recall:} say the Tamil for the day before yesterday, then the Tamil for a second (of time), then say \emph{añcalakam} again
\end{itemize}

\begin{tcolorbox}[breakable,colback=teal!4,colframe=teal!35!black,title={Before you move on}]
What does \ta{அஞ்சலகம்} mean? (\textquotedblleft{}A post office\textquotedblright{}.) Say it once more.
\end{tcolorbox}
\section[\texorpdfstring{\ta{நகரம்} (nakaram)}{நகரம் (nakaram)}]{\ta{நகரம்} (nakaram) — a city}

\label{lesson:TA-C234-nakaram}



\begin{quote}
Before the new one: say the Tamil for a second (of time), then the Tamil for a post office.
\end{quote}

\subsection*{You'll want to know: \ta{நகரம்}}
\textbf{\ta{நகரம்}} — \emph{nakaram} — \textquotedblleft{}a city\textquotedblright{}.

\begin{cousinweb}
\textbf{\ta{நகரம்}} (\emph{nakaram}) is native Dravidian. DEDR 3568 lists Malayalam \emph{nakar}, \textquotedblleft{}town\textquotedblright{}. The entry compares Sanskrit \emph{nagara}, \textquotedblleft{}town, city\textquotedblright{}.
\end{cousinweb}

\begin{grammarlens}[title={What you've built}]
One more word for this chapter.
\end{grammarlens}

\subsection*{Your turn}
\begin{itemize}
\raggedright
  \item \emph{Say it:} \emph{nakaram}
  \item \emph{Say it:} \emph{nakaram}, once more
  \item \emph{Say it:} say \emph{nakaram}
  \item \emph{Recall:} say the Tamil for a second (of time), then the Tamil for a post office, then say \emph{nakaram} again
\end{itemize}

\begin{tcolorbox}[breakable,colback=teal!4,colframe=teal!35!black,title={Before you move on}]
What does \ta{நகரம்} mean? (\textquotedblleft{}A city\textquotedblright{}.) Say it once more.
\end{tcolorbox}
\section[\texorpdfstring{\ta{நாடு} (nāṭu)}{நாடு (nāṭu)}]{\ta{நாடு} (nāṭu) — a country}

\label{lesson:TA-C234-natu}



\begin{quote}
Before the new one: say the Tamil for a post office, then the Tamil for a city.
\end{quote}

\subsection*{You'll want to know: \ta{நாடு}}
\textbf{\ta{நாடு}} — \emph{nāṭu} — \textquotedblleft{}a country\textquotedblright{}.

\begin{cousinweb}
\textbf{\ta{நாடு}} (\emph{nāṭu}) is native Dravidian. DEDR 3638 lists Malayalam \emph{nāṭu}, Kannada \emph{nāḍu} and Telugu \emph{nāḍu}.
\end{cousinweb}

\begin{grammarlens}[title={What you've built}]
One more word for this chapter.
\end{grammarlens}

\subsection*{Your turn}
\begin{itemize}
\raggedright
  \item \emph{Say it:} \emph{nāṭu}
  \item \emph{Say it:} \emph{nāṭu}, once more
  \item \emph{Say it:} say \emph{nāṭu}
  \item \emph{Recall:} say the Tamil for a post office, then the Tamil for a city, then say \emph{nāṭu} again
\end{itemize}

\begin{tcolorbox}[breakable,colback=teal!4,colframe=teal!35!black,title={Before you move on}]
What does \ta{நாடு} mean? (\textquotedblleft{}A country\textquotedblright{}.) Say it once more.
\end{tcolorbox}
\section[\texorpdfstring{\ta{தலைநகரம்} (talainakaram)}{தலைநகரம் (talainakaram)}]{\ta{தலைநகரம்} (talainakaram) — a capital city}

\label{lesson:TA-C234-talainakaram}



\begin{quote}
Before the new one: say the Tamil for a city, then the Tamil for a country.
\end{quote}

\subsection*{You'll want to know: \ta{தலைநகரம்}}
\textbf{\ta{தலைநகரம்}} — \emph{talainakaram} — \textquotedblleft{}a capital city\textquotedblright{}.

\begin{cousinweb}
\textbf{\ta{தலைநகரம்}} (\emph{talainakaram}) is \emph{talai}, \textquotedblleft{}head\textquotedblright{} (DEDR 3103), and \emph{nakaram}, \textquotedblleft{}city\textquotedblright{} (DEDR 3568): \textquotedblleft{}the head city\textquotedblright{}.
\end{cousinweb}

\begin{grammarlens}[title={What you've built}]
One more word for this chapter.
\end{grammarlens}

\subsection*{Your turn}
\begin{itemize}
\raggedright
  \item \emph{Say it:} \emph{talainakaram}
  \item \emph{Say it:} \emph{talainakaram}, once more
  \item \emph{Say it:} say \emph{talainakaram}
  \item \emph{Recall:} say the Tamil for a city, then the Tamil for a country, then say \emph{talainakaram} again
\end{itemize}

\begin{tcolorbox}[breakable,colback=teal!4,colframe=teal!35!black,title={Before you move on}]
What does \ta{தலைநகரம்} mean? (\textquotedblleft{}A capital city\textquotedblright{}.) Say it once more.
\end{tcolorbox}
\section[\texorpdfstring{\ta{நாற்சந்தி} (nāṟcanti)}{நாற்சந்தி (nāṟcanti)}]{\ta{நாற்சந்தி} (nāṟcanti) — a crossroads}

\label{lesson:TA-C234-narcanti}



\begin{quote}
Before the new one: say the Tamil for a country, then the Tamil for a capital city.
\end{quote}

\subsection*{You'll want to know: \ta{நாற்சந்தி}}
\textbf{\ta{நாற்சந்தி}} — \emph{nāṟcanti} — \textquotedblleft{}a crossroads\textquotedblright{}.

\begin{cousinweb}
\textbf{\ta{நாற்சந்தி}} (\emph{nāṟcanti}) is native \emph{nāl}, \textquotedblleft{}four\textquotedblright{} (DEDR 3655), and \emph{canti}, \textquotedblleft{}junction\textquotedblright{}, from Sanskrit \emph{saṁdhi} (CDIAL 12913): \textquotedblleft{}the four-way junction\textquotedblright{}.
\end{cousinweb}

\begin{grammarlens}[title={What you've built}]
One more word for this chapter.
\end{grammarlens}

\subsection*{Your turn}
\begin{itemize}
\raggedright
  \item \emph{Say it:} \emph{nāṟcanti}
  \item \emph{Say it:} \emph{nāṟcanti}, once more
  \item \emph{Say it:} say \emph{nāṟcanti}
  \item \emph{Recall:} say the Tamil for a country, then the Tamil for a capital city, then say \emph{nāṟcanti} again
\end{itemize}

\begin{tcolorbox}[breakable,colback=teal!4,colframe=teal!35!black,title={Before you move on}]
What does \ta{நாற்சந்தி} mean? (\textquotedblleft{}A crossroads\textquotedblright{}.) Say it once more.
\end{tcolorbox}
